\chapter{Magnetic resonance imaging (MRI)}
\index{Magnetic resonance imaging (MRI)}

\section*{Definition and Overview}
Magnetic resonance imaging (MRI) is a non-invasive, highly sophisticated diagnostic imaging modality that uses strong magnetic fields, radiofrequency pulses, and advanced computer algorithms to generate detailed, cross-sectional anatomical and physiological images of the human body. Unlike computed tomography (CT) or radiography, MRI does not employ ionising radiation, making it an exceptionally safe option for repeated examinations and paediatric populations. It offers superior soft-tissue contrast resolution, rendering it the diagnostic gold standard for evaluating the central nervous system, musculoskeletal system, cardiovascular system, and solid abdominal or pelvic organs.

\section*{Epidemiology}
Globally, the use of MRI has grown exponentially over the past few decades, with millions of scans performed annually. In many countries, public health-insurance statistics demonstrate a steady increase in MRI referrals, driven by expanding clinical indications, a growing ageing population, and increased availability of scanners. It is widely utilised across all age groups, though the highest volume of scans is performed on patients aged 40 to 75, primarily for musculoskeletal (e.g.\ knee and spine) and neurological indications. The availability of high-field magnets (1.5 Tesla and 3 Tesla) in regional and metropolitan areas has significantly enhanced patient access.

\section*{Aetiology and Pathophysiology}
The physiological principles of MRI rely on the magnetic properties of hydrogen nuclei (protons), which are abundant in human tissues due to the high water and fat content.
\begin{itemize}
    \item \textbf{Proton alignment:} Under normal conditions, hydrogen protons spin randomly. When the patient is placed in the strong static magnetic field of the scanner (B0), these protons align parallel or antiparallel to the field, creating a net longitudinal magnetisation vector.
    \item \textbf{Radiofrequency excitation:} A radiofrequency (RF) pulse is applied perpendicular to B0 at the Larmor frequency, tipping the protons into a transverse plane of magnetisation and putting them in phase (resonance).
    \item \textbf{Relaxation processes:} Once the RF pulse is turned off, protons return to their baseline state. This occurs via two pathways: longitudinal relaxation (T1 recovery, involving spin-lattice energy transfer) and transverse relaxation (T2 decay, involving spin-spin dephasing).
    \item \textbf{Signal detection and reconstruction:} The energy released as protons return to alignment induces an electrical current in receiver coils, which is converted via Fourier transform into detailed grey-scale images based on tissue-specific T1 and T2 relaxation times.
    \item \textbf{Contrast agents:} Gadolinium-based contrast agents (GBCAs) are often administered intravenously to alter local relaxation times, enhancing the visibility of vascular structures, inflammation, perfusion defects, and tumours.
\end{itemize}

\section*{Clinical Features}
An MRI scan is characterised by its multiplanar capability and unique patient experience, which can be challenging for some individuals.
\begin{itemize}
    \item \textbf{Multiplanar imaging:} Images can be acquired in axial, sagittal, coronal, or oblique planes without requiring the patient to change position.
    \item \textbf{Acoustic noise:} The rapid switching of gradient coils produces loud, repetitive tapping, clicking, or knocking sounds, necessitating the use of hearing protection (earplugs or headphones).
    \item \textbf{Spatial confinement:} The examination takes place within a long, narrow bore (cylinder), which can trigger claustrophobia or anxiety in susceptible individuals.
    \item \textbf{Immobility requirement:} The patient must remain completely still for the duration of each sequence (typically 2 to 7 minutes per sequence, with a total scan time of 15 to 45 minutes) to prevent motion artifacts.
    \item \textbf{Thermal effects:} The absorption of RF energy can cause a mild, generalised warming sensation in the tissues.
\end{itemize}

\section*{Differential Diagnosis}
While MRI is an investigation, choosing it over other imaging modalities requires understanding its relative strengths and limitations compared to alternative techniques.
\begin{itemize}
    \item \textbf{Computed tomography (CT):} CT is preferred for acute trauma, emergency evaluation of intracranial haemorrhage, and detailed cortical bone assessment, as it is much faster and more compatible with life-support equipment.
    \item \textbf{Ultrasonography:} Ultrasound is the first-line investigation for biliary disease (e.g.\ cholecystitis), pelvic pathology in females, and superficial soft tissues, as it is cheaper, dynamic, and does not require confinement.
    \item \textbf{Conventional radiography (X-ray):} Radiography remains the initial screening tool for suspected bone fractures or pulmonary pathology, providing rapid but limited two-dimensional structural detail.
    \item \textbf{Positron emission tomography (PET):} PET is a metabolic imaging modality used primarily in oncology to detect active cellular metabolism, often combined with CT or MRI (PET-CT/PET-MRI) for structural localisation.
\end{itemize}

\section*{When to Seek Medical Care}
Patients should consult their referring medical practitioner to discuss the indications, risks, and preparation required for an MRI.

\textbf{Contact your local emergency services for:}
\begin{itemize}
    \item Acute, life-threatening symptoms such as sudden severe headache, focal neurological deficits (signs of stroke), or unstable spinal injuries where rapid diagnosis via CT or urgent specialist assessment is required.
    \item Severe allergic reactions (anaphylaxis) occurring immediately after the injection of a gadolinium contrast medium, characterised by dyspnoea, facial swelling, or cardiovascular collapse.
\end{itemize}

\section*{Diagnosis and Investigations}
Before scheduling or performing an MRI, strict safety screening and preparation protocols must be followed to ensure patient safety.
\begin{itemize}
    \item \textbf{Safety screening questionnaire:} A mandatory, detailed checklist to identify metallic implants, pacemakers, cochlear implants, vascular clips, or metallic foreign bodies (e.g.\ intraocular metal shards) that could be dislodged or malfunction.
    \item \textbf{Renal function assessment:} Evaluating the estimated glomerular filtration rate (eGFR) before administering gadolinium-based contrast, due to the rare risk of nephrogenic systemic fibrosis (NSF) in patients with severe renal impairment.
    \item \textbf{Pregnancy screening:} While MRI is not known to harm the foetus, scans are generally avoided during the first trimester unless clinically essential and alternative imaging is unsuitable.
    \item \textbf{Anxiety and claustrophobia management:} Pre-treatment options include oral sedatives, utilising wide-bore or open scanners, or performing the procedure under general anaesthesia (particularly for paediatric patients).
\end{itemize}

\section*{Management}
The management of a patient undergoing an MRI involves patient preparation, monitoring, and post-procedure care.
\begin{itemize}
    \item \textbf{Patient preparation:} Removing all metallic objects, clothing with metal zippers, jewellery, hairpins, and transdermal patches containing metal foil.
    \item \textbf{Acoustic protection:} Providing earplugs or noise-cancelling headphones playing music to protect hearing and reduce distress.
    \item \textbf{Real-time monitoring:} Maintaining visual contact and intercom communication with the patient throughout the scan, and providing a hand-held squeeze ball (emergency alarm).
    \item \textbf{Post-contrast care:} Advising the patient to drink plenty of fluids to assist in clearing gadolinium from their system and monitoring the injection site for extravasation or delayed allergic reactions.
    \item \textbf{Sedation recovery:} Ensuring appropriate supervision and transport arrangements if oral or intravenous sedation was administered.
\end{itemize}

\section*{Complications}
Although MRI is safe, complications can arise from the physical environment, contrast media, or implanted devices.
\begin{itemize}
    \item \textbf{Projectile effect:} Ferromagnetic objects brought into the scanner room can become dangerous projectiles, drawn rapidly toward the center of the magnet.
    \item \textbf{Thermal burns:} Radiofrequency-induced currents in metallic implants, tattoos containing metallic ink, or touching wires (e.g.\ ECG leads) can cause localised skin burns.
    \item \textbf{Contrast reactions:} Mild side effects of gadolinium contrast include nausea, headache, and local injection site discomfort; anaphylactoid reactions are extremely rare.
    \item \textbf{Nephrogenic systemic fibrosis (NSF):} A rare, debilitating, systemic fibrosing disorder of the skin and internal organs associated with gadolinium use in patients with advanced renal failure.
    \item \textbf{Device malfunction:} The strong magnetic field can deactivate, reprogramme, or physically damage cardiac pacemakers, implantable cardioverter-defibrillators (ICDs), and neurostimulators.
\end{itemize}

\section*{Prognosis and Outcomes}
The diagnostic yield of MRI is exceptionally high, providing definitive anatomical detail that frequently guides clinical, surgical, or therapeutic management plans. The procedure itself has no long-term side effects when performed in accordance with safety guidelines. Following the scan, a specialist radiologist interprets the images, and the clinical report is forwarded to the referring doctor within a few days. The prognosis of the underlying condition depends on the specific pathology identified.

\section*{Prevention and Self-Care}
Patients can take several steps to ensure a smooth, comfortable, and safe MRI experience.
\begin{itemize}
    \item \textbf{Disclose medical history:} Inform the imaging staff of any metal implants, shrapnel, surgeries, pregnancy, or history of kidney disease.
    \item \textbf{Dress appropriately:} Wear comfortable, loose-fitting clothing without metal snaps, zippers, or underwires, or change into a hospital gown as requested.
    \item \textbf{Practice relaxation techniques:} Use deep breathing, visualisation, or counting strategies to remain calm and still inside the scanner.
    \item \textbf{Discuss claustrophobia early:} Request a mild sedative from your referring doctor beforehand if you anticipate severe anxiety or difficulty lying still.
    \item \textbf{Follow hydration instructions:} If contrast is administered, drink extra water following the procedure to support renal clearance of the contrast agent.
\end{itemize}

\section*{Sources and Further Reading}
The following reputable organisations provide detailed information regarding MRI procedures, safety guidelines, and clinical applications:
\begin{itemize}
    \item InsideRadiology (Royal Australian and New Zealand College of Radiologists). \textit{Detailed information on MRI procedures, safety, and patient experiences.} https://www.insideradiology.com.au/
    \item Healthdirect Australia. \textit{Consumer guide to magnetic resonance imaging, including preparation and what to expect.} https://www.healthdirect.gov.au/mri-scan
    \item Royal College of Radiologists. \textit{Clinical imaging guidelines and safety standards for MRI practice.} https://www.rcr.ac.uk/
    \item Mayo Clinic. \textit{Patient-focused overview of MRI scans, indications, risks, and procedures.} https://www.mayoclinic.org/tests-procedures/mri/about/pac-20384768
    \item NHS. \textit{Comprehensive guide to magnetic resonance imaging, including safety checks and uses.} https://www.nhs.uk/conditions/mri-scan/
\end{itemize}
